\section{The \textsc{CarbonShift} Algorithm and Guarantees}
\label{sec:algorithm}

\subsection{Algorithm}
\textsc{CarbonShift} is an online deferral algorithm that, at the arrival of each elastic job $j$, either admits the job immediately at a warm site whose current marginal intensity is below a slack-dependent threshold, or defers it into a short horizon of length $W\le\sigma_j$ during which a renewable-surplus window is forecast. The algorithm maintains, for every site $s$, a running estimate of the per-request embodied cost $\beta_s = e_s + E_s/R_s$, where $R_s$ is the cumulative number of requests served by warm instances on $s$; the embodied cost is updated whenever a job is admitted warm. The algorithm computes, for each candidate placement $(s,t)$ in the deferred set, a cost quote
\begin{equation}
q_j(s,t) \;=\; p_s\,\rho_j\,\tau_j\,\hat\mu_s(t) \;+\; \kappa_j(s,t)\,\beta_s,
\label{eq:quote}
\end{equation}
where $\hat\mu_s(t)$ is the marginal-intensity forecast (or, for the current slot, the observation) and $\kappa_j(s,t)$ is $0$ if a warm instance is projected to be available at $(s,t)$ and $1$ otherwise. The algorithm chooses the placement that minimises the quote, provided the quote is below a threshold $\theta_j$ that depends on the slack; otherwise it admits immediately at the cheapest current site. The pseudocode is given in Algorithm~\ref{alg:carbonshift}.

\begin{algorithm}[tb]
\caption{\textsc{CarbonShift} per-job decision}
\label{alg:carbonshift}
\begin{algorithmic}[1]
\Require job $j=(t^a_j,\tau_j,d_j,\mathcal{S}_j,\rho_j)$; current $\mu_s(t^a_j)$; warm set $W$; embodied cost $\beta_s$ for $s\in\mathcal{S}_j$; horizon $H$.
\If{$\sigma_j = d_j-\tau_j \le \sigma_{\min}$} \Comment{slack too small to defer}
  \State admit at $(s^*,t^a_j)$ minimising $q_j$ among $\mathcal{S}_j$ \textbf{return}
\EndIf
\State $\Theta_j \gets \{(s,t): s\in\mathcal{S}_j,\, t\in[t^a_j,\,t^a_j+\min(H,\sigma_j)]\}$
\State $\hat\mu_s(\cdot)\gets$ short-horizon forecast over $[t^a_j,t^a_j+\min(H,\sigma_j)]$
\For{$(s,t)\in\Theta_j$}
  \State $\kappa_j(s,t)\gets 1$ if no warm instance of $f$ projected on $s$ at $t$, else $0$
  \State $q_j(s,t)\gets p_s\rho_j\tau_j\hat\mu_s(t)+\kappa_j(s,t)\beta_s$
\EndFor
\State $\theta_j \gets \alpha\,\bar\mu_j + \beta^{\max}_j$ \Comment{threshold, see Theorem~\ref{thm:cr}}
\State $q^{*}\gets \min_{(s,t)\in\Theta_j} q_j(s,t)$
\If{$q^{*}\le \theta_j$}
  \State commit $(s^{*},t^{*}) = \arg\min q_j$
\Else
  \State admit at the current-slot site minimising $q_j$ at $t=t^a_j$
\EndIf
\State update $\beta_s$ for the chosen site
\end{algorithmic}
\end{algorithm}

\subsection{Assumptions and Scope}
We use three assumptions that match the observed behaviour of marginal-intensity series. (A1) The marginal-intensity series $\mu_s(t)$ at any site $s$ lies in $[\mu_{\min},\mu_{\max}]$ with $\Lambda = \mu_{\max}/\mu_{\min}$. (A2) The per-slot variation $|\mu_s(t+1)-\mu_s(t)|$ is bounded by $\Delta\mu$. (A3) The forecast $\hat\mu_s(t)$ over a horizon $H$ has bounded error $|\hat\mu_s(t)-\mu_s(t)|\le\eta$ for all $t$ in the horizon. (A1) and (A2) are empirical~\cite{chen2024towards,sukprasert2024implications}; (A3) is the standard short-horizon forecasting assumption.

\subsection{Competitive ratio}
We give the competitive guarantee in two regimes.

\begin{theorem}[Adversarial competitive ratio]
\label{thm:cr}
Under (A1), for arbitrary marginal-intensity sequences (no bounded-variation assumption), \textsc{CarbonShift} with $H=\infty$ and a perfect forecast has competitive ratio
\[
\rho \;\le\; 2\Lambda + 1,
\]
where $\Lambda = \mu_{\max}/\mu_{\min}$. The bound is tight up to a constant: there exists an arrival sequence and an intensity sequence achieving $\rho\ge\Lambda/4$.
\end{theorem}

\begin{theorem}[Bounded-variation competitive ratio]
\label{thm:cr-bv}
Under (A1) and (A2), \textsc{CarbonShift} with horizon $H\ge 1$ has competitive ratio
\[
\rho \;\le\; 1 + \frac{\Delta\mu\,H}{\mu_{\min}} + \frac{\beta^{\max}}{\mu_{\min}\tau_{\min}},
\]
where $\tau_{\min}$ is the minimum job duration and $\beta^{\max}$ the maximum embodied cost. Under (A3), replace $\Delta\mu$ by $\Delta\mu+2\eta$.
\end{theorem}

Theorem~\ref{thm:cr-bv} is the practically relevant bound: in real marginal traces, $\Delta\mu$ is an order of magnitude smaller than $\mu_{\max}$ over $H=12$ slots~\cite{chen2024towards}, so the bound is a small constant rather than $\Lambda$.

\subsection{Slack safety}
Deferral is not free: a deferred job consumes slack and risks SLA violation if the renewable-surplus window does not arrive. We define $\sigma_{\min}$ as the minimum slack below which deferral is disabled. The choice of $\sigma_{\min}$ trades carbon savings against SLA risk.

\begin{theorem}[Slack safety]
\label{thm:slack}
Under (A1), (A3) with forecast error $\eta$, the probability that \textsc{CarbonShift} defers a job with slack $\sigma_j$ and the deferred execution violates its deadline is bounded by
\[
\mathbb{P}[\text{SLA violation}\mid\sigma_j] \;\le\; \exp\!\Big(-\frac{\sigma_j-\sigma_{\min}}{H}\Big) + \mathbb{P}[\eta\text{ exceeds}\;\sigma_j-\tau_j],
\]
and the threshold $\sigma_{\min} = H+\tau_{\max}$ makes this probability zero for forecast errors within the (A3) envelope.
\end{theorem}

The theorem gives the orchestrator a knob: $\sigma_{\min}$ is the slack below which the controller conservatively admits. The knob trades carbon savings (more slack $\Rightarrow$ more deferral $\Rightarrow$ more savings) against SLA risk.

\subsection{Non-gameability}
A workload that misrepresents its deadline to get a longer deferral window could in principle reduce its carbon by claiming $\sigma_j>\sigma_j^{\text{true}}$. We show this does not help under \textsc{CarbonShift}.

\begin{theorem}[Non-gameability]
\label{thm:non-game}
Under (A1), a workload that reports a deadline $\tilde d_j > d_j$ (slack $\tilde\sigma_j>\sigma_j$) cannot reduce its charged carbon under \textsc{CarbonShift} compared with reporting the true deadline $d_j$.
\end{theorem}

The theorem matters because it makes the algorithm honest under a double-blind reviewer's favourite attack: a selfish tenant cannot game the carbon budget by misreporting its SLA.

\subsection{Algorithm summary}
Three properties follow. \textsc{CarbonShift} is competitive against an offline optimum that knows the full marginal trajectory (Theorem~\ref{thm:cr},~\ref{thm:cr-bv}); it is SLA-safe under a slack threshold that the operator can set (Theorem~\ref{thm:slack}); and it is non-gameable by deadline inflation (Theorem~\ref{thm:non-game}). The experiments in Section~\ref{sec:experiments} verify that the constants in Theorem~\ref{thm:cr-bv} are small in practice and that the SLA bound in Theorem~\ref{thm:slack} is conservative.
